% DONO: Nathan   ·   ESPAÇO: ~0,5 página (é a contribuição central)
\subsection{Verificação de conflitos}
% - Cada condição "qtd(p) op c" ou "total op c" vira uma faixa de inteiros
%   (domínio de intervalos; citar Cousot e Cousot \cite{cousot1977}).
% - Viabilidade: interseção das faixas + acoplamento
%   total = soma de qtd(p) * preço(p). Exemplo: qtd(camiseta) >= 3 implica
%   total >= R$ 149,70, incompatível com total < R$ 100,00.
% - Pior combinação: busca com/sem cada regra e duas podas (a inviabilidade
%   só cresce com o conjunto; limite pela soma restante). Exponencial no
%   pior caso: medir o tempo com 20, 40 e 80 regras.
% - Condições não analisáveis viram "pode valer sempre" (conservador): nunca
%   aceita um conflito real, mas pode recusar um programa correto.
% - Se fizer as tarefas (c04, m02): o que mudou no algoritmo.
